\uuid{w5n0}
\exo7id{1711}
\titre{exo7 1711}
\auteur{barraud}
\organisation{exo7}
\datecreate{2003-09-01}
\isIndication{false}
\isCorrection{false}
\chapitre{Réduction d'endomorphisme, polynôme annulateur}
\sousChapitre{Applications}
\cpge{mpsi-mp2i.calcul-matriciel-systemes.c, pcsi.calcul-matriciel-systemes-lineaires.c, ptsi.calcul-matriciel-systemes-lineaires.c, mpsi-mp2i.determinants.d, pcsi.calcul-matriciel-systemes-lineaires.d, ptsi.calcul-matriciel-systemes-lineaires.d, mpsi-mp2i.determinants.e, pcsi.determinants.d, mp-mpi.reduction-endomorphismes.c, pc.reduction.b, psi.reduction-endomorphismes.b, pt.reduction-endomorphismes.b}
\module{Algèbre}
\niveau{L2}
\difficulte{}

\contenu{
\question{
Soit $a$ un réel. On considère le système à $n$ équations et $n$
  inconnues suivant~:
  $$
  \left\{
    \begin{array}{lccccc}
               a\,x_{1}  - x_{2}  = 0 \\
      -x_{p-1}+a\,x_{p} -x_{p+1}  = 0 &  (2\leq p\leq n-1)\\
      -x_{n-1}+a\,x_{n}           = 0
    \end{array}
  \right.
  $$
  \'Ecrire la matrice $A_{n}$ associée à ce système. On note $D_{n}=\det
  A_{n}$. Calculer $D_{n}$ en fonction de $D_{n-1}$ et $D_{n-2}$
}
}
